\section{Zero avoidance costs information}
\label{sec:information}

The sampling and coverage results hold for every self-avoiding walk with
step bound \(D\). We now assume that the walk lies in
\(\mathcal A(\mathcal P_M)\), where \(\mathcal P_M\) is the union of
the prime batches in \eqref{eq:batches}. Thus it avoids zero modulo
both factors of every selected rational prime. This restriction forces
short words of increments to reveal information about the starting
residues. The purpose of this section is to quantify that information
at the common terminal time \(t_*\).

For each batch \(T\), choose independently a uniform subset of
\(q_T\) rational primes, with independent fair signs. Denote all
these auxiliary choices by \(\mathcal C\). They are independent of
the walk-sampling experiment. Entropies below are always computed with
\(\mathcal C\) fixed; an expectation \(\E_{\mathcal C}\) averages
the resulting entropies and does not include the choices as part of
the random data.

For a nonnegative integer time \(t\), write
\[
 W_L(t)=(z_{t+1}-z_t,\ldots,z_{t+L}-z_{t+L-1})
\]
for the next \(L\) increments. Each letter has at most \(K_D\)
possible values. Use the dyadic lengths
\[
 L_T=2^{\lfloor\log_2(T^{2/5})\rfloor},
 \qquad \tfrac12T^{2/5}<L_T\le T^{2/5}.
\]
These lengths are short enough that a word cannot hit one residue twice
modulo a factor of norm \(p\in[T,2T]\), yet long enough that the
displacement from \(A_0(T)\) has negligible entropy compared with the
information in the word. Their dyadic form will let us partition longer
words into shorter ones in Section~\ref{sec:budget}.

Let \(X'_T\) be the selected residue vector of batch \(T\) at
\(z_{t_*}\), and let \(O'_T\) contain the selected residues of all
smaller batches at that same point.

\begin{proposition}[Information cost of one batch]
\label{prop:information-rate}
There is a constant \(c_1>0\), depending only on the fixed schedule
parameters, such that for all sufficiently large \(M\) and every
batch in \eqref{eq:batches},
\begin{equation}
 \frac{1}{L_T}\,
 \E_{\mathcal C}
 I\bigl(X'_T;W_{L_T}(t_*)\mid O'_T\bigr)
 \ge \frac{c_1}{\log T}.
 \label{eq:information-rate}
\end{equation}
The threshold for \(M\) is uniform over all infinite self-avoiding
\(D\)-step walks in \(\mathcal A(\mathcal P_M)\).
\end{proposition}

\begin{proof}
Fix a batch, and abbreviate \(L=L_T\), \(q=q_T\), and
\(A_0=A_0(T)\). Let \(X\) be this batch's selected residue vector
at \(z_{A_0}\), and let \(O\) contain the selected residues of all
smaller batches at \(z_{A_0}\). Their counterparts at \(z_{t_*}\)
are \(X',O'\). Put
\[
 s=z_{t_*}-z_{A_0},\qquad W=W_L(t_*).
\]
We will first show that repeated observations of \((s,W)\) reveal a
fixed fraction of the entropy of \(X\) left after \(O\) is known.
Then we will charge the displacement entropy separately and move the
result to \((X',O',W)\).

By \eqref{eq:global-entropy} and \eqref{eq:smaller-alphabet},
\begin{equation}
 \E_{\mathcal C}H(X\mid O)\ge(1-2e_0)qL_*.
 \label{eq:initial-conditional-entropy}
\end{equation}
Indeed, \(H(X\mid O)\ge H(X)-H(O)\), and the logarithm of the
alphabet size of \(O\) is at most
\(e_0c_0T\le e_0qL_*\). The displacement estimate
\eqref{eq:checkpoint-displacement} gives
\begin{equation}
 H(s)=O_D(T^{r_0})=O_D(T^{1/4}).
 \label{eq:package-shift-entropy}
\end{equation}
Here \(s\) has the original schedule law, which is independent of
\(\mathcal C\).

\paragraph{Testing possible checkpoint residues.}
Choose a sufficiently large fixed constant \(K_0\), specified below,
and set
\[
 m_T=\left\lceil K_0(T/L)\log T\right\rceil.
\]
For fixed \(\mathcal C\), first draw \((O,X)\) with its original
law. Conditional on those values, draw \(m_T\) independent copies
\(P_a=(s_a,W^{(a)})\) of \((s,W)\); write
\(P=(P_1,\ldots,P_{m_T})\). A conditional law at a value of
\((O,X)\) of probability zero may be chosen arbitrarily, since it
does not affect this experiment. Each one-package marginal, together
with \((O,X)\), is the original schedule law.

For a selected factor \(\pi\) of norm \(p\), a candidate
\(x\in\Z[i]/(\pi)\) passes package \(a\) if
\[
 x+s_a+\sum_{h=0}^{l-1}W^{(a)}_h\not\equiv0\pmod\pi
 \qquad(0\le l<L),
\]
where words are indexed from zero and the empty sum is zero.
Let \(\mathcal L_\pi(P)\) be the candidates that pass every
package. The true coordinate \(X_\pi\) always belongs to this
list: the package's exact start has the prescribed residue \(X_\pi\),
and its subsequent positions lie on the zero-avoiding walk. Thus,
almost surely, these lists are nonempty, and
\begin{equation}
 H(X\mid O,P)
 \le \E\sum_{\pi\text{ selected}}\log|\mathcal L_\pi(P)|.
 \label{eq:passing-list-entropy}
\end{equation}
For each realized \((O,P)\), the conditional support of \(X\) is
contained in the Cartesian product of these lists. No independence
between the coordinates of \(X\) is needed.

\paragraph{Coverage after conditioning on residue data.}
The packages just constructed are independent given \((O,X)\),
whereas the continuations used in Proposition~\ref{prop:coverage}
were sampled from an exact checkpoint time. To connect these two
experiments, fix \(\mathcal C\) and a value of \((O,X)\) of
positive probability. Its conditional package law is a mixture over
the posterior exact times \(a\) at \(A_0\). At each such exact
time the future law is the original continuation law: \((O,X)\)
is determined by the time and hence imposes no further condition on
the fresh continuation randomness.
More explicitly, if $K_a$ is the fresh continuation law of $(s,W)$
from the exact time $a$, then for every set $E$ of packages,
\begin{equation}\label{eq:posterior-kernel}
 \Pr\{(s,W)\in E\mid O,X,\mathcal C\}
 =\sum_a\Pr\{A_0=a\mid O,X,\mathcal C\}\,K_a(E).
\end{equation}
The kernel $K_a$ does not depend on the auxiliary prime choices.

Call a selected factor \(\pi\) \emph{favorable} for these residue
data if the posterior probability that \eqref{eq:coverage} holds at
the exact \(A_0\) start is at least \(1/2\). Let \(\chi(a,\pi)\)
indicate failure of coverage for \(\pi\) at the exact start \(a\).
An unfavorable coordinate has posterior bad probability greater than
\(1/2\), so its indicator is at most twice that probability. The
tower property gives
\begin{align}
 \E_{\mathcal C}\E_{O,X}
 \#\{\text{unfavorable selected coordinates}\}
 &\le 2\E_{\mathcal C}\E_{O,X}
       \sum_{\pi\text{ selected}}
       \E[\chi(A_0,\pi)\mid O,X,\mathcal C] \notag\\
 &=2\E_{A_0}\E_{\mathcal C}
       \sum_{\pi\text{ selected}}\chi(A_0,\pi)
 \le 2\eta_0q.
 \label{eq:unfavorable-fraction}
\end{align}
The prior law of \(A_0\) is independent of \(\mathcal C\); its
posterior law need not be. At each exact time,
Proposition~\ref{prop:coverage} bounds the fraction of bad signed
primes by \(\eta_0\), and a uniform signed subset of size \(q\)
therefore contains at most \(\eta_0q\) bad coordinates on average.
This proves the last inequality.

Fix a favorable coordinate and let \(\mu\) be the posterior law
of its exact start. Write \(G\) for the starts where coverage
holds, so \(\mu(G)\ge1/2\). At every \(a\in G\), translating
the terminal residue law by \(z_a\bmod\pi\) gives
\[
 \Pr\{s\equiv-x\pmod\pi\mid A_0=a\}\ge\tau_0/p
\]
for all but fewer than \(p^{1-\beta_0}\) candidates \(x\). Define
\[
 u_x=\mu\left\{a\in G:
     \Pr\{s\equiv-x\pmod\pi\mid A_0=a\}\ge\tau_0/p\right\}.
\]
Averaging the exceptional-set sizes yields
\[
 \sum_{x\in\Z[i]/(\pi)}(\mu(G)-u_x)<p^{1-\beta_0}.
\]
If \(u_x<1/4\), the summand exceeds \(1/4\). Consequently all
but at most \(4p^{1-\beta_0}\) candidates satisfy \(u_x\ge1/4\).
We next show that each of these candidates is unlikely to survive
all the packages.

\paragraph{A short word supplies disjoint chances of rejection.}
For a candidate with \(u_x\ge1/4\), take an exact start counted
by \(u_x\). The smoothing estimate \eqref{eq:smoothing}, applied
at that start, shows for every \(0\le l<L\) that
\[
 \Pr\{z_{t_*+l}-z_{A_0}\equiv-x\pmod\pi\mid A_0=a\}
 \ge\tau_0/(2p).
\]
Indeed \(l\le T_+\), and \(T_+^{-9}\le\tau_0/(2p)\) for all
batches once \(M\) is sufficiently large. Mixing over the posterior
starts gives hit probability at least \(\tau_0/(8p)\) at each
offset.

For a single package, these \(L\) hit events are disjoint. If two
offsets hit, the two walk positions have equal residues modulo
\(\pi\). Their difference is nonzero by self-avoidance, and hence
has length at least \(\sqrt p\). Its length is also at most
\(DL<\sqrt p\), since \(L\le T^{2/5}\) and \(p\ge T\).
This contradiction establishes disjointness. One conditional package
therefore rejects the candidate with probability at least
\(cL/p\), where \(c=\tau_0/8\).

Independence of the packages given \((O,X)\) now yields, for every
favorable coordinate,
\begin{equation}
 \E[|\mathcal L_\pi(P)|\mid O,X]
 \le4p^{1-\beta_0}+p\exp(-cLm_T/p)
 \le5p^{1-\beta_0}.
 \label{eq:passing-list-size}
\end{equation}
Choose the fixed \(K_0\) large enough for the second inequality;
this is possible because \(T\le p\le2T\) and
\(Lm_T/p\ge(K_0/2)\log T\). The choices and the lower threshold
for \(M\) are uniform in the residue data and the walk.

By Jensen's inequality, the expected logarithm of a favorable list's
size is at most \((1-\beta_0)\log p+\log5\). For an unfavorable
coordinate use \(\log p\). Equations
\eqref{eq:passing-list-entropy}--\eqref{eq:passing-list-size} imply
\[
 \begin{aligned}
 \E_{\mathcal C}H(X\mid O,P)
 &\le(1-\beta_0)qL_*
       +2\eta_0\beta_0q\log(2T)+q\log5\\
 &\le(1-\beta_0/2)qL_*
 \end{aligned}
\]
for large \(M\). Here \(\eta_0=1/100\),
\(L_*\ge\log T\), and \(\log5/L_*\to0\).
Subtracting this bound from \eqref{eq:initial-conditional-entropy}
and using \(e_0=\beta_0/50\), we obtain
\begin{equation}
 \E_{\mathcal C}I(X;P\mid O)\ge(\beta_0/3)qL_*.
 \label{eq:package-information-lower}
\end{equation}
Repeated tests have thus exposed a fixed fraction of the checkpoint
residue entropy. It remains to place the cost in the terminal word
itself, since only terminal words will enter the common budget.

\paragraph{Transporting the information to the terminal time.}
For fixed \(\mathcal C\), conditional independence gives
\[
 \begin{aligned}
 I(X;P\mid O)
 &=H(P\mid O)-\sum_{a=1}^{m_T}H(P_a\mid O,X)\\
 &\le\sum_{a=1}^{m_T}I(X;P_a\mid O)
 =m_TI(X;s,W\mid O).
 \end{aligned}
\]
There are two costs in removing the displacement. First,
\(I(X;s\mid O)\le H(s)\). Second, given the full Gaussian-integer
displacement \(s\), residue addition is a bijection between
\((O,X)\) and \((O',X')\), so
\[
 I(X;W\mid O,s)=I(X';W\mid O',s).
\]
For any finite variables,
\[
 I(A;B\mid C,S)-I(A;B\mid C)
 =I(A;S\mid B,C)-I(A;S\mid C)\le H(S).
\]
Removing the conditioning on \(s\) therefore costs at most another
\(H(s)\). In total,
\begin{equation}
 I(X;s,W\mid O)\le I(X';W\mid O')+2H(s).
 \label{eq:translate-information}
\end{equation}
The one-package marginal is the original schedule law. Thus the
variables on the right are precisely
\(X'_T,O'_T,W_{L_T}(t_*)\), under the same terminal law as for
every other batch.

Finally, \(qL_*\ge c_0T\) and
\(m_T\le2K_0(T/L)\log T\) for large \(T\). Equations
\eqref{eq:package-information-lower},
\eqref{eq:translate-information}, and
\eqref{eq:package-shift-entropy} give
\[
 \E_{\mathcal C}I(X';W\mid O')
 \ge\frac{\beta_0c_0}{6K_0}\frac{L}{\log T}-O_D(T^{1/4}).
\]
Since \(T^{1/4}=o(L/\log T)\), this proves
\eqref{eq:information-rate} with, for example,
\(c_1=\beta_0c_0/(12K_0)\).
\end{proof}

\section{The common-law entropy budget}
\label{sec:budget}

Each batch now charges a positive amount of information to an increment
word. We must show that these charges can be added, despite the different
word lengths. The residue vectors will be nested in increasing batch
order. Knowing the earlier increments updates the residue vector to the
start of each new block, while the smoothing of \(t_*\) compares that
block's law with the law at \(t_*\). These two facts make the entropy
differences telescope.

\begin{lemma}[A common-law entropy telescope]
\label{lem:telescope}
For every fixed collection \(\mathcal C\) of prime choices,
\begin{equation}
 \sum_T\frac{I(X'_T;W_{L_T}(t_*)\mid O'_T)}{L_T}
 \le\log K_D+o(1),
 \label{eq:information-telescope}
\end{equation}
where the error tends to zero uniformly over the choices and the walks
as \(M\to\infty\).
\end{lemma}

\begin{proof}
Order the batches as \(T_1<\cdots<T_{n_B}\), and put
\(L_j=L_{T_j}\). For every time \(t\), let \(Q_j(t)\) contain
all selected residues from batches through \(T_j\) at \(z_t\);
let \(Q_0\) be empty. Define
\[
 a_j=\frac{H(W_{L_j}(t_*)\mid Q_{j-1}(t_*))}{L_j},
 \qquad
 b_j=\frac{H(W_{L_j}(t_*)\mid Q_j(t_*))}{L_j}.
\]
The \(j\)-th charge is \(a_j-b_j\). We claim that the initial
entropy rate for batch \(j+1\) is bounded, up to a small error,
by the remaining rate \(b_j\) after batch \(j\).

Both lengths are powers of two and \(L_j\le L_{j+1}\), so
\(L_j\) divides \(L_{j+1}\). Partition the longer word into
blocks of length \(L_j\). For a block starting at offset \(a\),
\(Q_j(t_*)\) and the preceding increments determine
\(Q_j(t_*+a)\), by adding their sum in every residue coordinate.
The entropy chain rule, followed by removal of conditioning, gives
\[
 H(W_{L_{j+1}}(t_*)\mid Q_j(t_*))
 \le\sum_{\substack{0\le a<L_{j+1}\\L_j\mid a}}
 H(W_{L_j}(t_*+a)\mid Q_j(t_*+a)).
\]

Both the block and its starting residue vector are deterministic
functions of the starting time on the fixed walk. By
\eqref{eq:smoothing}, their joint law at \(t_*+a\) differs from
that at \(t_*\) in total variation by at most \(T_+^{-9}\), since
\(a<L_{j+1}\le T_+^{2/5}\le T_+\). This is approximate invariance
of the sampled time law; no stationarity of the walk is assumed.

The logarithm of the full residue alphabet is at most
\[
 \sum_T k_T\log(2T)=O(T_+),
\]
by \eqref{eq:prime-count} and geometric spacing of the batches.
The block alphabet contributes at most
\(L_j\log K_D=O_D(T_+)\). Lemma~\ref{lem:continuity} therefore
gives, uniformly in the block and in \(\mathcal C\),
\[
 \begin{aligned}
 H(W_{L_j}(t_*+a)\mid Q_j(t_*+a))
 &\le H(W_{L_j}(t_*)\mid Q_j(t_*))+\epsilon_M,\\
 \epsilon_M&=O_D(T_+^{-7}).
 \end{aligned}
\]
Indeed the continuity bound is
\(2h(T_+^{-9})+O_D(T_+^{-8})\), which is smaller than the displayed
error for sufficiently large \(M\). Dividing the block sum by
\(L_{j+1}\) now yields
\[
 a_{j+1}\le b_j+\epsilon_M/L_j\le b_j+\epsilon_M.
\]
There are \(n_B=O(\log T_+)\) batches, so
\[
 \sum_{j=1}^{n_B}(a_j-b_j)
 \le a_1-b_{n_B}+(n_B-1)\epsilon_M
 \le\log K_D+o(1).
\]
Here \(b_{n_B}\ge0\), and the \(K_D\)-letter increment alphabet
gives \(a_1\le\log K_D\).
\end{proof}

This telescope is methodologically related to Tao's entropy-decrement
argument~\cite[Section~3, Equation~(40) and Lemma~17]{TaoEntropyDecrement}:
approximate translation invariance and conditional block subadditivity
limit accumulated information. Here the information charges arise from
zero avoidance across nested prime batches, and the preceding proof
establishes the required bound for our single terminal-time law.

\begin{proof}[Proof of Theorem~\ref{thm:sieve}]
Fix \(D\) and all schedule constants. Suppose that, for a given
sufficiently large \(M\), an infinite self-avoiding \(D\)-step walk
lies in \(\mathcal A(\mathcal P_M)\). Averaging
Lemma~\ref{lem:telescope} over \(\mathcal C\) and applying
Proposition~\ref{prop:information-rate} gives
\begin{equation}
 c_1\sum_T\frac1{\log T}\le\log K_D+o(1).
 \label{eq:final-information-bound}
\end{equation}

For each \(\lceil M/2\rceil\le m\le M\), the interval
\([X_m,1.05X_m]\) contains at least \(cX_m/\log B\) of the
numbers \(\log T=j\log B\), for an absolute \(c>0\) and all
sufficiently large \(M\). Its contribution to the sum in
\eqref{eq:final-information-bound} is at least
\(c/(1.05\log B)\). There are order \(M\) such disjoint
intervals. Their total contribution tends to infinity, whereas
\(c_1\), \(B\), and \(K_D\) are fixed. This contradicts
\eqref{eq:final-information-bound} once \(M\) is sufficiently large.

Every estimate and every lower threshold for \(M\) is uniform over
the hypothetical avoiding walks. Choose one such \(M\), depending
only on \(D\), and set \(\mathcal P_D=\mathcal P_M\). This is
one finite set of rational primes, with both factors imposed in
\(\mathcal A(\mathcal P_D)\), for which no infinite avoiding walk
exists. The auxiliary signed subsets were used only to average the
information inequalities; they do not make the final sieve depend on
the walk. This proves Theorem~\ref{thm:sieve} and supplies the input
to the periodicity argument in Section~\ref{sec:consequences}.
\end{proof}
